\documentclass[11pt]{article}
\usepackage[T1]{fontenc}
\usepackage[utf8]{inputenc}
\usepackage{lmodern}
\usepackage{amsmath,amssymb,amsthm}
\usepackage{microtype}
\usepackage{geometry}
\usepackage{enumitem}

\newtheorem{theorem}{Theorem}
\newtheorem{proposition}{Proposition}
\newtheorem{lemma}{Lemma}
\newtheorem{corollary}{Corollary}
\newtheorem{conjecture}{Conjecture}
\newtheorem{remark}{Remark}

\newtheorem{defn}[theorem]{Definition}
\newcommand{\VCdim}{\operatorname{VCdim}}

\newcommand{\cm}{\operatorname{cm}}


\title{A note on Open AI's "A counterexample to Hadwiger’s conjecture"}
\author{Pei Wu}
\date{\today}

\begin{document}
	
	\maketitle
	
	\section{Summary}
	In the note they construct a triangle-free graph $G$ that has many $4$-cycles, such that for the complement graph we has its independence number $\alpha(\bar G)=2$ and  connected matching number $\cm (\bar G)< m/100$ where $m$ is number of vertices. The connected matching number is defined as follows:
	
	\begin{defn}
		A matching $M$ in a graph $H$ is called a \emph{connected matching} if every two distinct edges of $M$ are joined by an edge of $H$: for all distinct $e,f\in M$, there exist $u\in V(e)$ and $v\in V(f)$ such that $uv\in E(H)$. The \emph{connected matching number} is
	\[
	\cm(H)=\max\bigl\{|M|: M\text{ is a connected matching in }H\bigr\}.
	\]
	\end{defn}
	
	The complement graph provides a counterexample to Hadwiger's conjecture, which asserts that $\chi(H)\le h(H)$, where $h(H)$ is the largest integer $t$ such that $H$ contains a $K_t$ minor. Indeed, Proposition~3.5 of the paper (proved in Appendix~\ref{app:minor-bound}) gives
	\[
	h(H)\le \frac{|V(H)|+4\cm(H)+2}{3}.
	\]
	Applying this to $H=\bar G$, and using $\alpha(\bar G)\le 2$ so that every color class has at most two vertices, we obtain, for $m$ large enough,
	\[
	h(\bar G)
	\le \frac{m+4\cm(\bar G)+2}{3}
	< \frac{26m}{75}+\frac{2}{3}
	< \frac{m}{2}
	\le \chi(\bar G).
	\]
	Thus $h(\bar G)<\chi(\bar G)$, contrary to Hadwiger's conjecture.
\appendix
\section{The clique-minor bound (Proposition 3.5)}
\label{app:minor-bound}

\begin{proposition}
For every finite nonempty graph $H$ with $m$ vertices,
\[
h(H)\le \frac{m+4\cm(H)+2}{3}.
\]
\end{proposition}

\begin{proof}
Write $c=\cm(H),\,b=h(H)$. Thus that there are $b$ pairwise disjoint, nonempty connected vertex sets, with an edge joining every pair of sets.
These sets are the branch sets of a $K_b$ minor. Let $s$ be the number of sets with one vertex and $e$ be the number of sets with two vertices.

Firstly $c>e+\frac{s-1}{2}$:
Each two-vertex branch set has an edge. And pairs of singleton branches supports an edge, the set of them is a connected matching, so we have the desired inequality. 

Then every remaining branch set has at least three vertices. Since the branch sets are disjoint,
\[
m\ge s+2e+3(b-s-e)=3b-2s-e.
\]
Combining two inequality give use the contradiction.
\end{proof}
\end{document}
